\chapter{Systemanalyse und Treiberentwurf}

\section{Analyse der Systemanforderungen}
Die Integration eines leistungsfähigen Grafiktreibers in den hhuOS-Mikrokernel
bringt im Vergleich zu monolithischen Umgebungen wie Linux eine Reihe
besonderer Herausforderungen mit sich. In hhuOS besitzt der Treiber keinen
inhärenten privilegierten Zugriff auf den Systemspeicher; stattdessen muss er
Ressourcen explizit von den Speicher- und PCI-Diensten des Kernels anfordern.

\subsection{Funktionale Anforderungen: Die Befehlspipeline}
Der Treiber muss einen bidirektionalen Kommunikationskanal zwischen dem
Gast-Betriebssystem und der Host-VMM bereitstellen.
\begin{enumerate}[label=\arabic*.]
    \item \textbf{PCI-Bus-Enumeration:} Der Treiber muss mit dem hhuOS-PCI-Manager interagieren, um das VirtIO-GPU-Gerät (ID \texttt{0x1050}) zu lokalisieren und seine herstellerspezifischen Capabilities zu analysieren, damit die MMIO-Offsets für die Konfiguration bestimmt werden können.
    \item \textbf{2D-Ressourcenverwaltung:} Der Treiber muss die Logik implementieren, um hostseitige Texturen (Resources) zu erstellen, Gast-RAM-Seiten als Backing Stores anzuhängen und den Lebenszyklus dieser Ressourcen zu verwalten, um Speicherlecks im Hypervisor zu verhindern.
    \item \textbf{Synchrone und asynchrone E/A:} Während Frame-Aktualisierungen (\textit{Flushes}) asynchron ablaufen können, erfordern Steuerbefehle, etwa das Erstellen einer Ressource, einen synchronen Handshake, um sicherzustellen, dass die Ressourcen-ID vor der weiteren Verwendung gültig ist.
\end{enumerate}

\subsection{Nicht-funktionale Anforderungen: Sicherheit und Leistung}
\begin{enumerate}[label=\arabic*.]
    \item \textbf{Strikte Isolierung:} Gemäß der Mikrokernel-Philosophie \cite{Liedtke1995} arbeitet der Treiber in seinem eigenen Adressraum. Fehler in der Behandlung des VirtIO-GPU-Protokolls müssen auf den Treiberprozess beschränkt bleiben, um eine systemweite Instabilität zu verhindern.
    \item \textbf{Interrupt-Effizienz:} In einer virtualisierten Umgebung sind Interrupts (VM-Exits) teuer. Der Entwurf muss, wo immer möglich, MSI-X nutzen und für Hochdurchsatz-Szenarien den Polling-Modus unterstützen, um den Overhead von Kontextwechseln zu minimieren.
\end{enumerate}

\section{Gesamtarchitektur des Treibers}
Der hhuOS-VirtIO-GPU-Treiber ist als modulares System konzipiert, das
hochrangige Grafikaufrufe von den Interaktionen mit der zugrunde liegenden
Hardware entkoppelt. Diese Modularität ist in Mikrokernel-Umgebungen, in denen
Dienstgrenzen strikt eingehalten werden, von entscheidender Bedeutung.

Die Architektur gliedert sich in drei primäre Funktionsbereiche:
\begin{enumerate}
    \item \textbf{Frontend-API-Schicht:} Diese Schicht stellt die Schnittstelle zum restlichen hhuOS-Ökosystem dar. Sie übersetzt generische Zeichen- oder Aktualisierungsanfragen in das logische VirtIO-GPU-Protokoll. Darüber hinaus verwaltet sie den virtuellen Framebuffer und stellt sicher, dass Pixeldaten auf Anwendungsebene korrekt für die GPU formatiert sind.
    \item \textbf{Zentrale VirtIO-Engine:} Dies ist das Herzstück des Treibers. Sie enthält die Protokollhandler, die \texttt{VIRTIO\_GPU\_CMD}-Pakete erzeugen. Außerdem verwaltet sie die beiden obligatorischen Virtqueues: die \textit{Control Queue} für Ressourcenverwaltung und Initialisierung sowie die \textit{Cursor Queue} für hardwarebeschleunigte Mauszeiger.
    \item \textbf{OS-spezifische HAL- und Schnittstellenschicht:} Diese Schicht bildet die Schnittstelle zu den hhuOS-Kerndiensten. Sie übernimmt das Einblenden von PCI-BARs in den Adressraum des Treibers, die Allokation physisch zusammenhängender DMA-Puffer sowie die Benachrichtigung des Hypervisors über die Notify-Register.
\end{enumerate}

\begin{center}
\begin{tikzpicture}[
    node distance=1.5cm,
    >=Stealth,
    comp/.style={rectangle, draw, text width=5cm, text centered, minimum height=1cm, rounded corners}
]
    \node[comp, fill=red!10] (app) {Benutzeranwendung (GUI/Kommandozeile)};
    \node[comp, below=of app, fill=orange!10] (api) {hhuOS-Grafik-API};

    \node[comp, fill=green!10, below=1.8cm of api] (core) {Zentrale VirtIO-GPU-Logik\\(Protokollhandler)};
    \node[comp, fill=green!10, below=0.8cm of core] (vq) {Virtqueue-Engine};

    \node[draw, dashed, inner sep=0.5cm, fit=(core) (vq), label=right:{\textbf{Treiberdomäne}}] (driverbox) {};

    \node[comp, fill=gray!20, below=1.2cm of vq] (hal) {OS-spezifische HAL};
    \node[comp, fill=yellow!10, below=of hal] (hw) {Hypervisor (QEMU/KVM)};

    \draw[->, thick] (app) -- (api);
    \draw[->, thick] (api) -- (core) node[midway, right] {Befehle};
    \draw[->, thick] (core) -- (vq);
    \draw[->, thick] (vq) -- (hal);
    \draw[->, thick, dashed] (hal) -- (hw) node[midway, right] {DMA/MMIO};
\end{tikzpicture}
\captionof{figure}{Mehrschichtige Treiberarchitektur in hhuOS}
\end{center}

\section{Zentrale Datenstrukturen}
Die Effizienz des Treibers hängt von Datenstrukturen ab, die präzise mit der
VirtIO-1.2-Spezifikation übereinstimmen. Da das Gerät, also der Hypervisor,
diese Strukturen direkt über DMA ausliest, sind Speicherausrichtung und
Padding von zentraler Bedeutung.

\subsection{Der Gerätezustand (VirtioGpuDevice)}
Die zentrale Klasse \texttt{VirtioGpuDevice} verwaltet den Laufzeitzustand der
Hardwareverbindung. Sie speichert die MMIO-Zeiger auf die PCI-BARs sowie die
Instanzen der Virtqueues.

\begin{lstlisting}[language=C++, basicstyle=\small\ttfamily, frame=single, caption=Zentrale Gerätestruktur]
class VirtioGpuDevice {
    // MMIO-Register
    volatile CommonCfg* common_cfg;
    volatile uint32_t* notify_reg;
    volatile GpuConfig* gpu_cfg;

    // Virtqueues
    Virtqueue control_vq;
    Virtqueue cursor_vq;

    // Ressourcenverwaltung
    uint32_t current_resource_id;
    DmaBuffer framebuffer;
};
\end{lstlisting}

\subsection{Die Split-Virtqueue-Strukturen}
Die \textbf{Split Virtqueue} besteht aus drei Speicherbereichen, die zwischen
Gast und Host gemeinsam genutzt werden. In hhuOS werden diese als ein einziger
zusammenhängender DMA-Block alloziert, um die Adressübersetzung zu
vereinfachen.

\begin{enumerate}
    \item \textbf{Deskriptortabelle:} Ein Array von 16-Byte-Einträgen, das Pufferadresse, Länge und Flags (z.\,B. \texttt{NEXT}, \texttt{WRITE}) definiert.
    \item \textbf{Available Ring:} Ein vom Gast beschriebener Ringpuffer, der Indizes von Deskriptoren enthält, die zur Verarbeitung bereitstehen.
    \item \textbf{Used Ring:} Ein vom Host beschriebener Ringpuffer, in dem das Gerät abgeschlossene Aufgaben meldet.
\end{enumerate}

\subsection{GPU-Steuerpakete}
Die gesamte Kommunikation über die Control Queue verwendet einen
Standard-Header (\texttt{virtio\_gpu\_ctrl\_hdr}). Diese Struktur stellt
sicher, dass sich Gast und Host über den Befehlstyp und eventuelle Fehlercodes
einig sind.

\begin{lstlisting}[language=C++, basicstyle=\small\ttfamily, frame=single]
struct virtio_gpu_ctrl_hdr {
    uint32_t type;     // VIRTIO_GPU_CMD_*
    uint32_t flags;
    uint64_t fence_id;
    uint32_t ctx_id;
    uint32_t padding;
};
\end{lstlisting}

\section{Detaillierter Initialisierungsablauf}
Die Initialisierung eines VirtIO-1.2-Geräts ist eine empfindliche
Zustandsmaschinen-Transition. Der hhuOS-Treiber stellt sicher, dass keine E/A
versucht wird, bevor der Host das \texttt{FEATURES\_OK}-Bit bestätigt hat.

\subsection{Phase 1: Erkennung und Capability-Mapping}
Im Gegensatz zu Legacy-PCI-Geräten, die feste BAR-Offsets verwenden, nutzt
VirtIO-GPU eine dynamische Liste von Capabilities. Der Treiber durchläuft die
verkettete Liste der PCI-Capabilities, um
\texttt{VIRTIO\_PCI\_CAP\_COMMON\_CFG} und
\texttt{VIRTIO\_PCI\_CAP\_NOTIFY\_CFG} zu finden. Dadurch kann der Treiber auch
moderne Hypervisoren unterstützen, bei denen die Benachrichtigungsregister in
einem anderen BAR liegen als die Konfigurationsregister.

\subsection{Phase 2: Der Handshake}
Der Treiber folgt der in \cite{Virtio12} definierten Sequenz:
\begin{enumerate}[label=\arabic*.]
    \item \textbf{Reset:} Das Statusregister wird gelöscht, um einen sauberen Ausgangszustand zu gewährleisten.
    \item \textbf{Feature-Verhandlung:} Der Treiber liest die 64-Bit-Feature-Bitmaske vom Gerät. Er muss \texttt{VIRTIO\_F\_VERSION\_1} aushandeln, um moderne 1.1-Features wie Little-Endian-Datenstrukturen zu aktivieren.
\end{enumerate}

\section{Entwurf der Hardwareabstraktionsschicht (HAL)}
Eine wichtige Entwurfsentscheidung war die Implementierung einer dünnen HAL.
Diese Schicht abstrahiert die kernelspezifischen Details von hhuOS, zum
Beispiel:
\begin{itemize}
    \item \textbf{Speicherabbildung:} Übersetzung physischer PCI-Adressen in den virtuellen Adressraum des Treibers über die Schnittstelle \texttt{os\_map\_mmio()}.
    \item \textbf{Physische Kontiguität:} Sicherstellung, dass DMA-Puffer physisch zusammenhängend sind. In hhuOS wird dies von einem speziellen DMA-Allokator übernommen, der die standardmäßige Heap-Allokation umgeht, um seitenausgerichtete und zusammenhängende Blöcke bereitzustellen.
\end{itemize}

\section{Speicherverwaltung: Die semantische Lücke}
Einer der komplexesten Aspekte der GPU-Virtualisierung ist die Überbrückung der
„semantischen Lücke“. Das Gast-Betriebssystem weiß, welche Pixel es zeichnen
möchte, aber die Host-GPU kann nicht direkt auf den privaten Speicher des Gasts
zugreifen.

\subsection{Direct Memory Access (DMA) und Pinnen}
Um dieses Problem zu lösen, implementiert der Treiber einen
Backing-Store-Mechanismus. Wenn ein Framebuffer in hhuOS alloziert wird, muss
der Treiber die entsprechenden Seiten pinnen. Dadurch wird sichergestellt, dass
der hhuOS-Kernel diese Seiten nicht auslagert oder neu zuweist, während die
Host-GPU daraus liest.

\begin{center}
\begin{tikzpicture}[>=Stealth, font=\small]
    \node[draw, fill=blue!10, text width=3.5cm, align=center] (gva) {Gast-virtuell\\(hhuOS-Anwendung)};
    \node[draw, fill=green!10, right=2cm of gva, text width=3.5cm, align=center] (gpa) {Gast-physisch\\(RAM)};
    \node[draw, fill=red!10, right=2cm of gpa, text width=3.5cm, align=center] (hpa) {Host-physisch\\(VRAM)};

    \draw[->, thick] (gva) -- (gpa) node[midway, above] {MMU};
    \draw[->, thick] (gpa) -- (hpa) node[midway, above] {EPT/SLAT};

    \node[below=0.5cm of gpa, text width=5cm, align=center, font=\itshape] {Der VirtIO-GPU-Treiber arbeitet auf dieser Ebene und übergibt GPA an den Host.};
\end{tikzpicture}
\captionof{figure}{Dreistufige Adressübersetzung bei VirtIO}
\end{center}

\section{Virtqueue-Entwurf und Synchronisation}
Die Virtqueue ist die primäre Abstraktion für asynchrone Datenübertragung in
VirtIO. Für den GPU-Treiber wird eine \textbf{Split-Virtqueue}-Architektur
verwendet. Dieser Entwurf basiert auf einem gemeinsam genutzten
Ringpufferprotokoll, das teure VM-Exits durch die Bündelung von Befehlen
minimiert.

\subsection{Deskriptorverkettung und Anordnung}
Jeder GPU-Befehl besteht aus einem Anforderungs-Header, der vom Gerät gelesen
wird, und einem Antwort-Header, der vom Gerät beschrieben wird. Da sich beide
an unterschiedlichen Speicherorten befinden, wird \textbf{Deskriptorverkettung}
verwendet. Ein einzelner logischer Befehl wird dabei durch zwei oder mehr
Deskriptoren dargestellt, die über das \texttt{NEXT}-Flag miteinander
verknüpft sind.

\begin{center}
\begin{tikzpicture}[node distance=0.5cm, font=\small]
    \node[draw, fill=blue!5, minimum width=4cm] (desc0) {Desc 0: [Header GPA, Länge, NEXT]};
    \node[draw, fill=green!5, below=of desc0, minimum width=4cm] (desc1) {Desc 1: [Antwort GPA, Länge, WRITE]};

    \draw[->, thick] (desc0.east) .. controls +(right:1cm) and +(right:1cm) .. (desc1.east)
        node[midway, right] {Verkettet};

    \node[draw, dashed, inner sep=0.2cm, fit=(desc0) (desc1), label=left:{Logischer Befehl}] {};
\end{tikzpicture}
\captionof{figure}{Virtqueue-Deskriptorverkettung für GPU-Befehle}
\end{center}

\subsection{Der Synchronisations-Handshake}
Um die Konsistenz zwischen dem hhuOS-Treiber als Produzent und dem Hypervisor
als Konsument zu wahren, wird ein strenges Speicherordnungssprotokoll
eingehalten. Dadurch wird verhindert, dass der Host teilweise geschriebene
Deskriptoren liest.

\begin{enumerate}[label=\arabic*.]
    \item \textbf{Payload schreiben:} Der Treiber schreibt die Befehlsdaten in den DMA-Puffer.
    \item \textbf{Deskriptor einrichten:} Der Treiber füllt die Deskriptortabelle mit der physischen Adresse des Puffers.
    \item \textbf{Speicherbarriere (SFENCE):} Ein \texttt{SFENCE}-Befehl stellt sicher, dass alle vorherigen Speicheroperationen für andere Prozessoren und die Host-CPU sichtbar sind, bevor der Available-Index aktualisiert wird \cite{Virtio12}.
    \item \textbf{Available-Index aktualisieren:} Der Treiber erhöht den Index im Available Ring.
    \item \textbf{Benachrichtigen (Kick):} Der Treiber schreibt in das PCI-Notify-Offset. Dies löst einen VM-Exit aus und signalisiert dem Host, mit der Verarbeitung zu beginnen.
\end{enumerate}

\subsection{Der Benachrichtigungszyklus im Überblick}
Die Interaktion zwischen der Benachrichtigung durch den Treiber und der
späteren Rückmeldung des Geräts wird im Folgenden dargestellt:

\begin{center}
\begin{tikzpicture}[
    node distance=2cm,
    >=Stealth,
    state/.style={
        rectangle,
        rounded corners,
        draw,
        text width=1.6cm,
        align=center,
        minimum height=1.2cm
    }
]
    \node[state, fill=blue!10] (start) {Treiber\\Vorbereitung};
    \node[state, fill=gray!10, right=of start] (avail) {Available\\Ring\\aktualisieren};
    \node[state, fill=red!10, right=of avail] (kick) {Türklingel\\(Notify)};
    \node[state, fill=green!10, below=of kick] (host) {Host\\Verarbeitung};
    \node[state, fill=blue!10, left=of host] (used) {Used Ring\\aktualisieren};

    \draw[->, thick] (start) -- (avail);
    \draw[->, thick] (avail) -- (kick);
    \draw[->, thick] (kick) -- node[right] {VM-Exit} (host);
    \draw[->, thick] (host) -- (used);
    \draw[->, thick] (used) -- node[left] {Interrupt} (start);
\end{tikzpicture}
\captionof{figure}{Asynchroner Befehlslebenszyklus und Signalisierung}
\end{center}

\paragraph{Konkretes Beispiel: \texttt{RESOURCE\_FLUSH} über die Control Queue}
Abbildung~\ref{fig:virtqueue-resource-flush} zeigt einen beispielhaften
Befehlsablauf über die Control Queue am Beispiel von
\texttt{RESOURCE\_FLUSH}. Der Treiber verbindet zunächst Anfrage und Antwort
über eine Deskriptorkette, trägt den Kopf der Kette in den Available Ring ein
und benachrichtigt den Host über das Notify-Register. Nach der Verarbeitung
aktualisiert der Host den Used Ring und stellt die Antwort im vorgesehenen
Puffer bereit. Der Treiber kann danach den Abschluss des Befehls auswerten und
die weitere Verarbeitung fortsetzen.

\begin{center}
\begin{tikzpicture}[
    node distance=1.6cm and 1.4cm,
    box/.style={
        draw,
        rounded corners=3pt,
        align=center,
        minimum width=3.2cm,
        minimum height=1.1cm,
        font=\small
    },
    arrow/.style={
        -{Latex[length=2.5mm]},
        thick
    },
    step/.style={
        circle,
        draw,
        inner sep=1pt,
        minimum size=16pt,
        font=\scriptsize
    }
]

\node[box] (req) {Request-Header\\Descriptor 0};
\node[box, right=of req] (resp) {Response-Puffer\\Descriptor 1};

\node[box, below=1.8cm of req] (avail) {Available Ring};
\node[box, right=of avail] (notify) {Notify-Register};
\node[box, right=1.8cm of notify] (host) {Host verarbeitet\\\texttt{RESOURCE\_FLUSH}};
\node[box, below=1.8cm of host] (used) {Used Ring};
\node[box, left=of used] (driver) {Treiber liest\\Antwort aus};

\draw[arrow] (req) -- (resp);
\node[step, above=1mm of $(req)!0.5!(resp)$] {1};

\draw[arrow] (req) -- (avail);
\node[step, left=2mm of $(req)!0.5!(avail)$] {2};

\draw[arrow] (resp) -- (avail);
\node[step, right=2mm of $(resp)!0.5!(avail)$] {2};

\draw[arrow] (avail) -- (notify);
\node[step, above=1mm of $(avail)!0.5!(notify)$] {3};

\draw[arrow] (notify) -- (host);
\node[step, above=1mm of $(notify)!0.5!(host)$] {4};

\draw[arrow] (host) -- (used);
\node[step, right=1mm of $(host)!0.5!(used)$] {5};

\draw[arrow] (used) -- (driver);
\node[step, above=1mm of $(used)!0.5!(driver)$] {6};

\end{tikzpicture}

\vspace{0.5em}

\begin{minipage}{0.9\textwidth}
\footnotesize
\textbf{Ablauf:}
(1) Anfrage und Antwort werden über eine Deskriptorkette verbunden.
(2) Der Kettenkopf wird in den Available Ring eingetragen.
(3) Der Treiber benachrichtigt den Host über das Notify-Register.
(4) Der Host verarbeitet den Befehl \texttt{RESOURCE\_FLUSH}.
(5) Der Abschluss wird im Used Ring eingetragen.
(6) Der Treiber liest die Antwort aus.
\end{minipage}

\captionof{figure}{Beispielhafter Befehlsablauf über die Control Queue am Beispiel von \texttt{RESOURCE\_FLUSH}}
\label{fig:virtqueue-resource-flush}
\end{center}

\subsection{Synchrone Wait-Used-Logik}
Obwohl die Virtqueue grundsätzlich asynchron arbeitet, benötigt hhuOS für
Einrichtungsbefehle, z.\,B. \texttt{RESOURCE\_CREATE\_2D}, ein synchrones
Verhalten. Der Treiber implementiert daher eine Polling-Schleife, die den
\texttt{used\_idx} des Used Rings überprüft. Um Systemabstürze zu verhindern,
ist ein Watchdog-Timer integriert:
\begin{equation}
    T_{\text{verstrichen}} > T_{\text{timeout}} \implies \text{Geräte-Reset}
\end{equation}
Wenn der Host den Deskriptor nicht innerhalb von 1000\,ms zurückgibt,
protokolliert der hhuOS-Kernel einen fatalen Treiberfehler und versucht, die
VirtIO-Zustandsmaschine zurückzusetzen.

\section{Fazit des Systementwurfs}
Die in diesem Kapitel durchgeführte Analyse bestätigt, dass das
VirtIO-1.2-Protokoll in Kombination mit einer robusten HAL und einer strengen
Speicherverwaltung in hhuOS die notwendige Infrastruktur für einen
leistungsfähigen Grafikstack bereitstellt. Durch die Betonung der
Paravirtualisierung gegenüber der Emulation wird der VM-Exit-Overhead erheblich
reduziert und die Grundlage für die in Kapitel~4 beschriebene Implementierung
geschaffen.

\newpage